\section{Conclusion}\label{sec:conclusion}

The formalization proves a specific chain of results.  Arbitrary-term
polynomial reduction is terminating over an arbitrary commutative coefficient
ring.  Under the unit-leading-coefficient hypothesis, reduction equivalence
modulo $G$ agrees with congruence modulo $\ideal{G}$, and the Gr\"obner
condition is equivalent to reduction-to-zero and the confluence properties.
Assuming additionally that the coefficient ring has no zero divisors,
Buchberger's S-polynomial condition is equivalent to the Gr\"obner condition.
On the state space carrying the unit-leading-coefficient invariant,
Noetherianity prevents infinite sequences of Buchberger-state steps.  One step
and finite chains preserve finiteness, so finite input yields a finite terminal
Gr\"obner basis in the finite-variable field specialization.  These theorems
provide a verified relational specification of the abstract Buchberger
procedure.  A practical certified solver additionally requires finite
executable data structures, a verified normal-form function, and a computable
strategy for selecting polynomial pairs and processing their S-polynomials.

\section*{Acknowledgments}
The author thanks the Lean and Mathlib communities for the existing library,
discussion, and review on which this work builds.

\section*{Funding}
This work was supported by the National Research Foundation of Korea (NRF)
grant funded by the Korea government (MSIT) (No. RS-2025-00520280).

\section*{Data and software availability}
No external dataset was generated or analyzed.  The canonical Lean source is
publicly available in the author's \lean{mathlib4} fork.  The reproducible
companion artifact, CI configuration, and generated documentation are
available in the \lean{Buchberger} repository and through the artifact homepage
\url{https://sanghyeok0.github.io/Buchberger/}.  The exact source and artifact
snapshots, Lean toolchain, and reproduction commands are recorded in
Section~\ref{sec:reproducibility}.
